% \usetikzlibrary {shapes.multipart}

% \def \EventTreeHeaderBackgroundFill {none}
% \def \EventTreeHeaderLabelBackgroundFill {black!5}
% \tikzset{
%     headlabel/.style={
%         % font=\tiny,
%     },
% }

% \newcommand{\HeaderNode}[2]{#1 \nodepart[headlabel]{two}#2}

% \usetikzlibrary{graphs,graphdrawing,trees}
% \usetikzlibrary{calc,positioning}

\begin{figure}[ht!]
\centering
\begin{tikzpicture}
% \tikzset{grow'=right,level distance=48pt}
% \tikzset{execute at begin node=\strut}
% \tikzset{every tree node/.style={anchor=base west}}
\tikzset{
    grow'=right,
    level distance=48pt,
    execute at begin node=\strut,
    every tree node/.style={anchor=base west},
    edge from parent/.append style={very thick},
    edge from parent/.style={
        draw,
        edge from parent path={
            (\tikzparentnode.east) -| ($(\tikzparentnode.east)!0.5!(\tikzchildnode.west)$) |- (\tikzchildnode.west)
        },
    },
    every node/.style={anchor=center,font=\small\bfseries, text centered, inner sep=0pt},
    every level 0 node/.style={circle, font=\small\bfseries, draw, fill=blue!30, inner sep=0pt},
    every internal node/.style={font=\small, inner sep=4pt},
    every leaf node/.style={rectangle, draw, fill=blue!30, minimum width=2.5cm, text centered},
    % frontier/.style={distance from root=400pt},
    frontier/.style={distance from root=\textwidth-20pt-100pt},
}
% \tikzset{
%     grow'=right,
%     level distance=50pt,
%     sibling distance=10pt,
%     execute at begin node=\strut,
%     every tree node/.style={anchor=base west},
%     edge from parent/.style={
%         draw,
%         edge from parent path={
%             (\tikzparentnode.east) -| ($(\tikzparentnode.east)!0.5!(\tikzchildnode.west)$) |- (\tikzchildnode.west)
%         },
%     },
%     every node/.style={
%         % anchor=base north,
%         font=\small,
%         text centered, 
%         inner sep=1pt,
%         % anchor=base west,
%     },
%     every level 0 node/.append style={
%         circle,
%         fill=blue!40,
%         % minimum size=4pt,
%         draw,
%         % xshift=-4pt,
%         % inner sep=0pt,
%         % outer sep=0pt,
%         % text depth=0pt,
%         % text width=0pt,
%         % text height=0pt,
%     },
%     % every internal node/.style={
%     %     inner xsep=0pt,
%     %     outer xsep=0pt,
%     %     minimum size=0pt,
%     %     text width=0pt,
%     % },
%     % every leaf node/.style={
%     %     rectangle, 
%     %     fill=black!5,
%     %     align=center,
%     %     text width=43pt,
%     % },
%     frontier/.style={distance from root=\textwidth-20pt-100pt},
% }
\Tree [.\(I\)
    [.\(F_1^{\text{succ}}\)
        [.\(F_2^{\text{succ}}\)
            [.\(X_1\) ]
        ]
        [.\(F_2^{\text{fail}}\)
            [.\(X_2\) ]
        ]
    ]
    [.\(F_1^{\text{fail}}\)
        [.\(X_3\) ]
    ]
]
\end{tikzpicture}
% \begin{tikzpicture}
% \begin{scope}
% \graph[
%         grow right=50pt,
%         nodes={
%             anchor=south,
%             minimum width=48pt,
%             minimum height=40pt,
%             text width=42pt,
%             text depth=28pt,
%             line width=0pt,
%             align=flush left,
%             outer sep=0pt,
%             inner xsep=4pt,
%             inner ysep=5pt,
%         },
%         edge={draw=none},
%         ] {
%   %     % Functional Events Header
%     {
%        [
%         nodes={
%             shape=rectangle split,
%             rectangle split parts=2,
%             rectangle split part fill={\EventTreeHeaderBackgroundFill, \EventTreeHeaderLabelBackgroundFill},
%             rectangle split ignore empty parts=true,
%             draw,
%             font=\small,
%             },
%         ]
%       "\HeaderNode{Initiating Event}{IE}" --
%       "\HeaderNode{Functional Event 1}{$F_1$}" --
%       "\HeaderNode{Functional Event 2}{$F_2$}" -- 
%       "\HeaderNode{Functional Event 3}{$F_3$}" --
%       "\HeaderNode{End State}{ES}"
%   };
% };
% \end{scope}
% \begin{scope}[yshift=-61pt]
% \tikzset{
%     grow'=right,
%     level distance=50pt,
%     sibling distance=0pt,
%     execute at begin node=\strut,
%     every tree node/.style={anchor=base west},
%     edge from parent/.style={
%         draw,
%         thick,
%         edge from parent path={
%             (\tikzparentnode.east) -| ($(\tikzparentnode.east)!0.5!(\tikzchildnode.west)$) |- (\tikzchildnode.west)
%         },
%     },
%     every node/.append style={
%         anchor=base west,
%     },
%     every level 0 node/.style={
%         circle,
%         fill=blue!40,
%         minimum size=4pt,
%         draw,
%         xshift=-4pt,
%         inner sep=0pt,
%         outer sep=0pt,
%         text depth=0pt,
%         text width=0pt,
%         text height=0pt,
%     },
%     every internal node/.style={
%         inner xsep=0pt,
%         outer xsep=0pt,
%         minimum size=0pt,
%         text width=0pt,
%     },
%     every leaf node/.style={
%         rectangle, 
%         fill=black!5,
%         align=center,
%         text width=43pt,
%     },
%     frontier/.style={distance from root=\textwidth-20pt-100pt},
% }
% \Tree [.{I} 
%         [.{} [.{} [.{}  {IE•LD•LI•FCA} ]
%                     [.{} {IE•LD•LI•FCA'} ]]
%               [.{} {IE•LD•LI'} ]]
%         [.{} [.{} {IE•LD'•LI} ]
%                [.{} [.{} {IE•LD'•LI'•FCA} ]
%                       [.{} [.{} {IE•LD'•LI'•FCA'•SEO} ] [.{} {IE•LD'•LI'•FCA'•SEO'} ] ] ] ] ]
% \end{scope}
% \end{tikzpicture}
\caption{Illustrative event tree with an initiating event \(I\), two functional events \(F_1\) and \(F_2\), and three end-states \(X_1, X_2, X_3\).}
\label{fig:event_tree_example}
\end{figure}